\section{Appendix: proofs of \cref{thm:soundness-v21,thm:saturation-v22}}\label{app:proof}

\subsection{Auxiliary results}

\begin{lemma}[Strict descent]\label{lem:descent}
Let \(s = (s_V, s_P) \in \mathcal{S}\) and let
\(N(s) \;\stackrel{\mathrm{def}}{=}\; \#\{R \in \Lambda : s \not\models_R\}\),
the number of layout regions that fail per-region consistency.
If \(s \neq \bot\) and \(N(s) > 0\), then either
\(\mathcal{G}(s) = \bot\) or \(N(\mathcal{G}(s)) < N(s)\).
\end{lemma}

\begin{proof}
By construction of \(\mathcal{G}\) (\cref{def:G}). Step 1
(decode) examines each inconsistent region. For each such region,
either every subblock RS-decodes within capacity (no residual
syndrome), in which case the region's syndrome becomes zero after
the decoded content is installed and \(N\) decreases by one for
that region; or at least one subblock yields a residual syndrome,
in which case step 5 emits the journal entry and the operator
returns \(\bot\) by I5. In either case, \(N\) cannot remain
unchanged: an iteration that does not return \(\bot\) strictly
decreases \(N\).
\end{proof}

\begin{proposition}[Fail-closed absorption]\label{prop:absorb}
If \(\mathcal{G}(s) = \bot\) for some \(s\), then
\(\mathcal{G}(\bot) = \bot\).
\end{proposition}

\begin{proof}
Direct from invariant I5: fail-closed monotonicity. The proposition
restates I5 in operator form.
\end{proof}

\subsection{Proof of \cref{thm:soundness-v21}}

\begin{proof}[Proof of \cref{thm:soundness-v21}]
Fix a consistent initial state \(s_0 \models\), an event
\(e \in \mathcal{E}_{\mathrm{EMR}}\), and let \(s_1 = e(s_0)\).
Define the iteration \(s_{i+1} = \mathcal{G}(s_i)\) for
\(i \geq 1\), with the convention that if \(s_i = \bot\) then
\(s_{i+1} = \bot\) (by \cref{prop:absorb}, this is consistent
with the operator's behaviour).

By the EMR-bounded hypothesis on \(\mathcal{E}_{\mathrm{EMR}}\),
no subblock of \(s_1\) carries more than \(t\) symbol errors. The
RS decoder is well-defined within capacity, returning the unique
codeword within Hamming distance~\(t\) of the received word
(MacWilliams--Sloane~\cite{macwilliams1977}, Ch.~9); the kernel
\texttt{lib/reed\_solomon}~\cite{kernelrslib} implements this. Hence step 1 of \cref{def:G} succeeds
on every region of \(\Lambda\) on which \(s_1\) differs from
\(s_0\), without invoking the saturation handling of step 5.

By \cref{lem:descent}, if \(s_i \neq \bot\) and \(N(s_i) > 0\),
then \(N(s_{i+1}) < N(s_i)\) and \(s_{i+1} \neq \bot\) (because
the decoder is within capacity by hypothesis). Since
\(N(s_i) \in \{0, 1, \dots, |\Lambda|\}\), there exists a least
\(K \le |\Lambda|\) such that either \(N(s_K) = 0\) or
\(s_K = \bot\). The first case implies \(s_K \models\) by
\cref{def:consistent}; the second is the fail-closed exit.

The fail-closed exit is reached only if a sentinel verification
(step 4) failed, which is a violation of I2 not covered by the
RS-decode hypothesis. Such a violation can occur if the EMR
event corrupted a region in a way that bypasses the RS protection
(for instance, by writing canonical-looking bytes whose CRC
happens to validate, an event of probability bounded by \(2^{-c}\) for a \(c\)-bit
CRC; BEAMFS v2 uses CRC-32, so this probability is approximately
\(2.3 \cdot 10^{-10}\) per region per RS-decode event); the
operator's response is the designed safe failure. This case is bounded by I2's sentinel separation
property and does not affect the termination claim.
\end{proof}

\subsection{Proof of \cref{thm:saturation-v22}}

\begin{proof}[Proof of \cref{thm:saturation-v22}]
Fix a consistent \(s_0\) and an event
\(e \in \mathcal{E}_{\mathrm{sat}}\). Let \(s_1 = e(s_0)\). By
hypothesis, there exists at least one subblock of \(s_1\) whose
symbol-error count exceeds \(t\). Apply \(\mathcal{G}\) to \(s_1\).

Step 1 of \cref{def:G} (decode) invokes the RS decoder on every
subblock of every inconsistent region. On the saturating
subblock, the decoder returns a non-zero residual syndrome (the
canonical out-of-capacity behaviour~\cite{kernelrslib}). The
operator detects the residual and proceeds to step 5.

Step 5 (saturation log) calls \texttt{beamfs\_log\_rs\_event}
with \texttt{positions} \(=\) \texttt{NULL},
\texttt{n\_positions} \(=\) 0, and
\texttt{flags |= BEAMFS\_RS\_EVENT\_FLAG\_UNCORRECTABLE}, with
the affected \((\text{block}, \text{subblock})\) coordinates
identified by the RS-decode failure. The journal entry is
committed to persistent storage by the standard buffer-write
sequence (\texttt{mark\_buffer\_dirty},
\texttt{sync\_dirty\_buffer}). We assume the implementation commits the journal entry to
persistent storage before returning, an assumption realized by the
kernel function \texttt{beamfs\_log\_rs\_event} which calls
\texttt{sync\_dirty\_buffer} before its return path; this is
verifiable by audit and is the on-disk format v4
contract~\cite{beamfs-impl}.

The operator returns \(\bot\). Conclusion (1) holds. The journal
entry committed in step 5 satisfies conclusion (2): it is tagged
with \texttt{UNCORRECTABLE}, and its \((\text{block},
\text{subblock})\) coordinates correspond to a saturating subblock
of \(e(s_0)\) by construction.
\end{proof}

\subsection{Remark on falsification scope}

Both proofs assume that the kernel implementation realizes the
operator definition of \cref{def:G} faithfully. Implementation
faithfulness is not the subject of the theorem and is addressed
by the traceability table of \cref{tab:m2c-trace}, by
\texttt{checkpatch.pl} compliance per \cref{sec:m2c}, and
ultimately by mechanized verification (v3+ work). A
counter-example to either theorem produced by RadFI or by an
equivalent fault-injection methodology against the actual kernel
binary may be either a counter-example to the theorem
(\(\mathcal{E}_{\mathrm{EMR}}\) is too narrow,
\(\mathcal{E}_{\mathrm{sat}}\) is mis-stated) or a counter-example
to the implementation faithfulness (a kernel bug, not a
mathematical defect). Distinguishing the two cases is the proper
diagnosis task following any falsification report.
